\section{Introduction}
\label{sec:intro}

A server converts most of its electrical input into heat. In a home, that heat can displace
heat-pump or backup-heater operation, while the server sells computing work. The appeal is simple:
one electricity bill may provide both a computing service and useful heat. The hard question is
whether the two uses coincide often enough, and whether the value created covers the equipment and
operation costs for both the operator and the host.

This question is especially sharp for a single-family house. In the Basel reference case, hourly
weather in 2025 implies a peak demand equivalent to 12.5 servers of 1~kW, while summer hot water
uses only about 0.6 servers. A system sized for the winter peak is therefore poorly utilised in
summer; a small system leaves much of the winter demand to backup heat. Scheduling can shift work
within those physical limits, but cannot create compute customers or heat demand. The business
case must consequently be tested against both the heat balance and the price of sellable work.

We study the following questions:
\begin{description}
  \item[Q1] What heat coverage, utilisation and deadline performance can a case-compliant single
    house achieve, first with unlimited work and then with finite, sourced job demand?
  \item[Q2] Under which conditions does the single-house system cover its hardware, and how many
    servers per house does that favour?
  \item[Q3] In a multi-site setting, can moving deadline-bound work to sites with year-round heat
    demand reduce missed deadlines?
\end{description}

\paragraph{Approach.} We use the case's fixed Basel reference house and hourly weather, and add
labelled variants: a house calibrated to an organiser-provided gas record, a local scenario with an
indoor pool and a synthetic Swiss network. An event-driven fleet model draws jobs from sourced
volumes and traces, quotes prices and schedules work with one heat-first policy shared by the
simulation and the real-container controller. A rolling-horizon linear programme and a full-year
upper bound serve as analytical comparisons. For economics, we separate the case break-even
(electricity net of heat value) from full cost after hardware, and expose the volume, price and
server count at which the house breaks even. Network results are scenario analyses, not forecasts
of a deployed market.

\paragraph{Contributions.} This paper makes four claims, each limited to the evidence available.
\begin{enumerate}
  \item A reproducible reference-house analysis that separates the physical limit (unlimited work)
    from the result under sourced finite demand, and locates the full-cost boundary in volume,
    price and server count (Sections~\ref{sec:case}, \ref{sec:economics} and~\ref{sec:evaluation}).
  \item One heat-first scheduling policy used by the fleet simulation and a real-container controller,
    with greedy, network-greedy and linear-programme strategies as research comparisons
    (Sections~\ref{sec:model} and~\ref{sec:scheduling}).
  \item A check of the case's house model against a real house's measured gas use, and network
    scenarios showing when routing to year-round heat sinks changes deadline outcomes
    (Section~\ref{sec:evaluation}).
  \item A demonstrator in which the same scheduler freezes and resumes a real transcription and
    language-model container on the heat signal, with measured CPU time and energy
    (Section~\ref{sec:demonstrator}).
\end{enumerate}

The paper's practical target is not to claim that any server makes a house cheaper to heat. It is
to identify a credible single-house path: enough paid work at a sufficient realised price, a
hardware and installation cost compatible with that revenue, and heat that displaces an actual
alternative. If those conditions cannot be met, the result should make that failure boundary
legible to a homeowner and to DigitalHeat.
